\section{Resumen}
Este documento presentará el diseño e implementación de un sistema híbrido de control y
automatización para una grúa portacontenedores portuaria de muelle tipo pórtico
(\textit{STS: Ship-to-Shore}), orientado a la operación semi-automática coordinada de sus dos
movimientos principales: la traslación del carro (\textit{Trolley}) y el izaje de la carga
(\textit{Hoist}). El proyecto integrará los conocimientos del espacio curricular Autómatas y
Control Discreto, combinando un control de eventos discretos (DEDS) con un control regulatorio
en tiempo discreto, modelados y simulados en Matlab/Simulink/Stateflow.

En primer lugar se construirá el modelo matemático de la planta física, conformada por los
subsistemas de carro e izaje, acoplados entre sí por la carga suspendida que se balancea como
un péndulo. El modelo contemplará la transmisión por cables de acero elásticos, los frenos
mecánicos, el perfil de obstáculos relevado por un sensor LIDAR, los sensores de
realimentación y los moduladores de torque de los accionamientos.

Sobre dicha planta se diseñará una arquitectura de control jerárquica de tres niveles: el
Nivel 2 (regulatorio) generará las consignas de torque mediante lazos de velocidad y un control
de balanceo de la carga; el Nivel 1 (supervisor) será un autómata que coordinará los modos de
operación, la lógica de toma y liberación de la carga y la planificación de trayectorias
seguras; y el Nivel 0 (protección) sobrescribirá a los niveles superiores y llevará la grúa a un
estado seguro ante cualquier anomalía.

Finalmente, el comportamiento del sistema se validará mediante simulación, incorporando una
animación 2D en tiempo real desarrollada en Python para visualizar la operación, y se trasladará
la lógica de control a un PLC virtual programado en CODESYS, comunicado con Simulink
mediante el protocolo industrial OPC-UA.


\newpage
\section{Introducción}
Este trabajo tiene como objetivo el control de los accionamientos eléctricos y la operación
semi-automática coordinada de una grúa portacontenedores STS,
empleada para la carga y descarga de contenedores en grandes barcos de transporte marítimo. Se
consideran los dos movimientos principales, continuos y simultáneos, de la grúa en un plano
vertical (2D): la traslación horizontal del carro y el izaje vertical de la carga bajo su
posición. Ambos movimientos se encuentran dinámicamente acoplados a través de la carga suspendida, que se comporta
como un péndulo, y deben respetar restricciones de posición, velocidad y aceleración.

El sistema físico se acciona mediante cables de acero impulsados desde mecanismos rotativos con
tambores, cajas de engranajes, frenos mecánicos y motores eléctricos, gobernados por
accionamientos de velocidad variable. La carga (\textit{headblock/spreader}, con o sin
contenedor) se modela como una masa concentrada con dos grados de libertad, cuya dinámica es
híbrida: cambia discretamente según el estado de los \textit{twistlocks} y según su interacción
de contacto con el perfil de obstáculos del entorno. La grúa dispone, además, de un conjunto de 
sensores de realimentación que incluyen encoders, finales de carrera, detectores de sobrevelocidad, 
una celda de carga, un sensor de ángulo de balanceo y un LIDAR para relevar el perfil de obstáculos. 
Estos dispositivos proporcionan la información necesaria para la operación y las funciones de protección.

Sobre esta planta física se implementa una arquitectura de control jerárquica compuesta por tres niveles: 
un nivel regulatorio encargado de los lazos de control discretos, un nivel supervisor responsable de la 
coordinación de las maniobras de manera segura y un nivel de protección dedicado a llevar al sistema a estado seguro
ante situaciones de emergencia.

El proyecto se organiza en distintas etapas: el modelado matemático de la planta física, el
diseño e implementación de la arquitectura de control, la validación del sistema mediante simulación con una animación 2D en tiempo real
y la traducción de la lógica de control a un PLC programado en CODESYS para su ejecución en un
entorno realista.

% \begin{figure}[ht]
%     \centering
%     \includegraphics[width=1\textwidth]{imagenes/sistema_fisico.png}
%     \caption{Robot manipulador elemental de 1 g.d.l. en plano vertical (péndulo rígido actuado)}
% \end{figure}
